\uuid{y6rE}
\exo7id{1194}
\titre{exo7 1194}
\auteur{legall}
\organisation{exo7}
\datecreate{1998-09-01}
\isIndication{false}
\isCorrection{false}
\chapitre{Suite}
\sousChapitre{Convergence}
\cpge{mpsi-mp2i.nombres-reels-suites.c, pcsi.nombres-reels-suites.b, ptsi.nombres-reels-suites.b, mpsi-mp2i.nombres-reels-suites.d, pcsi.nombres-reels-suites.c, ptsi.nombres-reels-suites.c, mpsi-mp2i.nombres-reels-suites.f, pcsi.nombres-reels-suites.e, ptsi.nombres-reels-suites.e, bcpst1.suites-reelles, mp-mpi.topologie-espaces-normes.b, pc.espaces-vectoriels-normes.b, psi.espaces-vectoriels-normes.b}
\module{Analyse}
\niveau{L1}
\difficulte{}

\contenu{
\question{
Soit $  (u_n)_{n\in { \Nn}}  $ une suite born\' ee de nombres r\' eels telle que
$  \displaystyle{\lim _{n\rightarrow \infty }(u_{n+1}-u_{n})=0}  .$ Montrer que les valeurs
d'adh\`erence
de la suite $  (u_n)_{n\in { \Nn}}  $ forment un intervalle de $  { \Rr}  .$
}
}
